\documentclass[11pt]{article}
\usepackage[margin=1in]{geometry}
\usepackage{amsmath,amssymb,amsthm}
\usepackage{booktabs}
\usepackage[hidelinks]{hyperref}
\usepackage{url}
\emergencystretch=2em

\newtheorem{theorem}{Theorem}
\newtheorem{corollary}[theorem]{Corollary}
\theoremstyle{definition}
\newtheorem{definition}{Definition}
\theoremstyle{remark}
\newtheorem{remark}{Remark}

\newcommand{\lean}[1]{\texttt{\detokenize{#1}}}

\title{Mathematics Without Human Understanders:\\
A Lean-Checked Model of Producer-Independence and Trust Transfer\\[4pt]
\large Draft}
\author{Vasily Shablov\\[2pt]
\small Formalization and drafting assisted by Claude (Anthropic)}
\date{October 2026}

\begin{document}
\maketitle

\begin{abstract}
We examine two claims about mathematics produced by artificial agents. The first is that an
artifact can be valid mathematics even if no human understands it, or could. The second is
that, given a long enough record of verified superiority, it becomes rational to defer to such
an agent more than to the combined human community, just as we already defer to doctors and
engineers. We build a small formal model containing a Hilbert-style (formalist) definition of
mathematics, a Thurston-style (humanist) definition parameterised by a predicate
\emph{IsHuman}, a checker-based account of trust, a two-type Bayesian model of
trust transfer, and a gradual-replacement account of the boundary of ``human''. Sixteen
results are proved in Lean~4 with Mathlib. They show: the formalist definition is
producer-independent; the humanist definition inherits every open question about the boundary
of \emph{IsHuman} and collapses to the formalist one as that boundary widens; trusting a
checked proof never requires trusting its producer; verified successes push trust above any
threshold $t<1$ but never to certainty; and no evidence from checkable tasks constrains
behaviour on uncheckable ones. Lean certifies only that the conclusions follow from the
assumptions. The philosophical content lies in the assumptions, and we state each one
explicitly.
\end{abstract}

\section{Introduction}

A recent joke depicts two people celebrating that an AI ``proved that the Riemann zeta function
is zero-free in the half-plane $\Re(s) > 7/8$'' and then admitting that neither knows what
$\Re(s)$ means. The joke lands because the celebration is right: a zero-free half-plane
$\Re(s) > \sigma$ for any fixed $\sigma<1$ is unknown (the ``quasi-Riemann hypothesis''), and
proving it would be a major result in analytic number theory. Yet neither person is in a
position to know that. They are correct by trust, not by understanding.

The joke points to two questions.
\begin{enumerate}
\item \textbf{Is it mathematics?} If an artifact is correct but no human understands it, or
even could, is it still mathematics?
\item \textbf{Should we believe it?} If an agent repeatedly proves things humans cannot, and
is rarely or never caught in an error, when, if ever, does its word outweigh that of human
experts?
\end{enumerate}
A common answer to the first question, following Thurston~\cite{thurston}, is that the goal of
mathematics is human understanding, so an incomprehensible proof is at best a degenerate case.
A common answer to the second is that AI excels only at tasks with binary, checkable answers,
while the questions that generate mathematics (``what \emph{is} geometry?'') have no reward
signal to train on.

We defend a different position.
\begin{itemize}
\item[(T1)] Validity is a property of the derivation, not of its producer or audience
(formalism, in the tradition of Hilbert~\cite{zach}).
\item[(T2)] Trust in a sufficiently reliable agent can rationally exceed trust in human
experts. The point at which this happens in practice is a ``fuzzy boundary'': a threshold set
partly by psychology and sociology rather than by logic alone.
\end{itemize}
A thesis of this kind cannot be ``proved'' outright. What can be proved is that it follows from
explicit assumptions. Making those assumptions precise is useful in itself: it shows where the
real disagreements lie. Section~\ref{sec:model} develops the model. Section~\ref{sec:discussion}
discusses what is and is not established.

\paragraph{Prior art.} The arguments below are not new in substance. The examples of accepted
mathematics that no human has checked line by line include the four colour
theorem~\cite{appelhaken}, the Kepler conjecture~\cite{flyspeck}, the Boolean Pythagorean
triples problem~\cite{bpt} and Keller's conjecture in dimension seven~\cite{keller}.
Psychological research on trust in algorithms includes~\cite{dietvorst}. Gradual-replacement
arguments in philosophy of mind include~\cite{chalmers}. Our contribution is a single
machine-checked model that puts these threads side by side.

\section{The model}\label{sec:model}

All results are in the Lean file \path{TrustThesis.lean}. Section~\ref{sec:lean} maps
theorem numbers to Lean names.

\subsection{Two definitions of mathematics}

\begin{definition}
A \emph{proof system} $S$ consists of a type of formulas, a type of proofs and a relation
$\mathrm{Proves}(\pi,\varphi)$. An \emph{artifact} is a pair (producer, proof).
\end{definition}

\begin{definition}[Formalist]
An artifact $a$ is \emph{formal mathematics} for $\varphi$ if $\mathrm{Proves}(a.\mathrm{proof},
\varphi)$.
\end{definition}

\begin{definition}[Humanist]
Given a predicate $\mathrm{IsHuman}$ on agents and a relation $U(h,\pi,\varphi)$ (``$h$
understands $\pi$ as a proof of $\varphi$''), $a$ is \emph{humanist mathematics} for $\varphi$
if it is formal mathematics and some $h$ with $\mathrm{IsHuman}(h)$ satisfies
$U(h, a.\mathrm{proof}, \varphi)$.
\end{definition}

\begin{theorem}[Producer irrelevance]
Formal mathematics does not depend on the producer: $(x,\pi)$ and $(y,\pi)$ are formal
mathematics for the same formulas.
\end{theorem}

\begin{theorem}
Humanist mathematics is formal mathematics.
\end{theorem}

\begin{theorem}[Monotonicity]
If $\mathrm{IsHuman}_1 \subseteq \mathrm{IsHuman}_2$, then humanist mathematics with respect to
$\mathrm{IsHuman}_1$ is humanist mathematics with respect to $\mathrm{IsHuman}_2$.
\end{theorem}

\begin{theorem}[Collapse]
If every agent counts and every producer understands its own valid proofs, humanist and formal
mathematics coincide.
\end{theorem}

\begin{theorem}[Boundary dependence]
Suppose a valid proof is understood by exactly one agent $s$, with
$\neg\mathrm{IsHuman}_1(s)$ and $\mathrm{IsHuman}_2(s)$. Then the proof is humanist mathematics
for $\mathrm{IsHuman}_2$ but not for $\mathrm{IsHuman}_1$.
\end{theorem}

Theorem~5 makes precise the observation that the humanist definition ``contains conjectures''.
Whether a synthetic brain, an uploaded mind or a language model counts as a possible
understander is an open question. The humanist definition makes the mathematical status of a
proof depend on that question. Theorems~3 and~4 show the direction of travel: the more agents
we admit as understanders, the closer the humanist definition comes to the formalist one.

\subsection{The trusted base of a checked proof}

\begin{definition}
A checker $c$ is \emph{sound} if $c(\pi,\varphi) = \mathrm{true}$ implies
$\mathrm{Proves}(\pi,\varphi)$.
\end{definition}

\begin{theorem}[Trusted base]
Let $c$ be a sound checker. Let the logic be sound with respect to an intended meaning:
$\mathrm{Proves}(\pi,\varphi) \Rightarrow \mathrm{Meaning}(\varphi)$. Let the statement be faithful:
$\mathrm{Meaning}(\varphi) \Leftrightarrow Q$. Then any artifact accepted by $c$ for
$\varphi$ establishes $Q$, whoever produced it.
\end{theorem}

The producer does not appear among the hypotheses. This is the asymmetry that makes mathematics
special: finding a proof may require superhuman ability, but checking one is mechanical, and a
proof kernel such as Lean's~\cite{lean4} is small enough to be audited. The intelligence gap
between the producer and the reader therefore matters less for proofs than for any other kind
of expert output.

Three things remain to be trusted: the checker, the logic and the faithfulness of the statement.
The third is the weak point, since a correct proof of a mis-stated theorem is worthless. One
proposal is to have other agents audit the statement.

\begin{theorem}[Independent audits]
If $n$ auditors miss an unfaithful statement independently, with probabilities
$\varepsilon_i \le e$, the probability that all miss it is at most $e^n$.
\end{theorem}

\begin{theorem}
If $0 \le e < 1$ and $\delta > 0$, then $e^n < \delta$ for all sufficiently large $n$.
\end{theorem}

The model assumes independence. This is the substantive assumption: auditors built on similar
data and methods tend to make correlated errors, and correlation breaks the $e^n$ bound.

\subsection{Evidential trust transfer}

Fix two hypotheses about an agent. With prior probability $p>0$ it is of a \emph{high} type,
solving each independent verified task correctly with probability $a_{\mathrm{hi}}$. Otherwise
it is of a \emph{low} type, with accuracy $a_{\mathrm{lo}} < a_{\mathrm{hi}}$. After $n$
consecutive verified successes the posterior probability of the high type is
\[
P_n = \frac{p\,a_{\mathrm{hi}}^n}{p\,a_{\mathrm{hi}}^n + (1-p)\,a_{\mathrm{lo}}^n}.
\]

\begin{theorem}
If $p>0$ and $0 \le a_{\mathrm{lo}} < a_{\mathrm{hi}}$, then $P_n \to 1$.
\end{theorem}

\begin{theorem}[The boundary is crossed]
For every threshold $t < 1$ there is an $N$ such that $P_n > t$ for all $n \ge N$.
\end{theorem}

\begin{theorem}[Never certainty]
If $0 < p < 1$ and $a_{\mathrm{lo}}, a_{\mathrm{hi}} > 0$, then $P_n < 1$ for every $n$.
\end{theorem}

\begin{theorem}[Outweighing a community]
Model the aggregated judgement of a human community as a fixed reliability $r$. If
$r < a_{\mathrm{hi}}$, the agent's posterior expected accuracy
$P_n a_{\mathrm{hi}} + (1-P_n)a_{\mathrm{lo}}$ exceeds $r$ for all sufficiently large $n$.
\end{theorem}

Theorems~10 and~11 together give a precise version of the ``fuzzy boundary''. Evidence alone
crosses any threshold short of certainty in finite time, but never reaches certainty itself. A
person who demands zero risk ($t=1$) is therefore never convinced by evidence, however long the
record. Where a person or a society actually sets $t$ is not determined by the mathematics. It
is a psychological and social parameter. Studies of \emph{algorithm aversion}~\cite{dietvorst}
find that people abandon an algorithm after seeing it err, even when it outperforms humans.
This amounts to setting $t$ much higher for machines than for people.

\emph{Numerical illustration.} Take a sceptical prior $p = 0.01$, $a_{\mathrm{hi}} = 0.999$ and
$a_{\mathrm{lo}} = 0.9$. The posterior first exceeds $0.5$, $0.9$, $0.99$ and $0.999$ after
$45$, $66$, $89$ and $111$ consecutive verified successes respectively. The expected accuracy
first exceeds a community reliability of $r=0.95$ after $45$ successes, and $r=0.99$ after
$67$.

\subsection{The limits of transfer}

Let tasks be split into \emph{checkable} and \emph{uncheckable}. An observer sees only an
agent's correctness on checkable tasks.

\begin{theorem}[Evidence is blind off the checkable set]
Any quantity computed from observations on checkable tasks takes the same value for two agents
that agree on checkable tasks.
\end{theorem}

\begin{corollary}
There is always an agent that is wrong on every uncheckable task yet produces exactly the same
evidence as a perfect agent.
\end{corollary}

So the conclusions of Section~2.3 hold on the checkable domain only. Extending them to
uncheckable domains requires an additional \emph{transfer assumption}: that accuracy carries
over, with no distribution shift and no deception. This assumption cannot be tested with
evidence of the kind that motivated the trust in the first place. Research on scalable
oversight, such as~\cite{debate}, aims to enlarge the checkable set, for instance by letting
agents check each other.

\subsection{What is a human?}

\begin{theorem}[Replacement dilemma]
Let $f(0), \dots, f(n)$ be a chain of agents, for instance a brain with $i$ neurons replaced by
functionally identical synthetic ones. If $\mathrm{IsHuman}(f(0))$ and
$\neg\mathrm{IsHuman}(f(n))$, then there is a single step $i<n$ with $\mathrm{IsHuman}(f(i))$
and $\neg\mathrm{IsHuman}(f(i+1))$.
\end{theorem}

\begin{theorem}[Functional duplicates]
If understanding depends only on input--output behaviour, a behavioural duplicate of a human
understands everything that human understands.
\end{theorem}

Theorem~15 is the formal core of Chalmers' gradual-replacement arguments~\cite{chalmers} and of
the Ship of Theseus. Anyone who denies that a 1:1 synthetic brain is human must accept that
replacing one particular neuron makes the difference. Theorem~16 shows that under a behavioural
notion of understanding, a 1:1 synthetic brain qualifies as an understander in Thurston's sense,
whatever its substrate.

Current language models are not functional duplicates of any human brain, so neither theorem
settles their status. Whether they belong to the class of understanders is, in the terms of
Theorem~5, exactly the kind of open question on which the humanist definition depends. We take
no position on it here.

\section{Discussion}\label{sec:discussion}

\paragraph{What Lean certifies.} Every theorem above is proved in Lean~4 with Mathlib and
depends on at most the standard axioms (\lean{propext}, \lean{Classical.choice},
\lean{Quot.sound}). The mathematics is elementary. The point of formalising it is not
difficulty but precision: the hypotheses of each theorem are the complete list of what the
corresponding philosophical claim needs. Anyone who rejects a conclusion must reject a
named hypothesis.

\paragraph{The hypotheses that carry the weight.}
\begin{itemize}
\item \emph{Faithfulness of statements} (Theorem~6). This is where humans, or auditors, remain
necessary. Auditing by other agents helps only to the extent that their errors are independent.
\item \emph{Fixed community reliability} (Theorem~12). By Condorcet's jury
theorem~\cite{condorcet}, a large community of independent judges with individual reliability
above $1/2$ has majority reliability approaching $1$. The claim that one agent can outweigh
``all mathematicians combined'' therefore relies on human errors being correlated, which is
plausible since mathematicians share training, methods and blind spots, but it is an
assumption.
\item \emph{Transfer} (Corollary~14). Trust earned on checkable tasks does not, by itself,
license trust on uncheckable ones.
\end{itemize}

\paragraph{On ``disproving G\"odel''.} The incompleteness theorems are themselves proved in
weak systems, and have been machine-checked~\cite{paulson}. A valid ``refutation'' could only
mean exhibiting an inconsistency in the foundations, that is, a derivation of a contradiction.
Such a derivation is a finite object and can be checked mechanically, however incomprehensible
its structure. It would be valid mathematics in the sense of (T1), and an extreme instance of
Theorem~6.

\paragraph{Machines and the ``why'' questions.} An objection to (T2) is that AI can only answer
well-posed yes/no questions, while mathematics grows from open questions with no reward signal.
The model shows that this matters less for whether to trust an AI's conclusions than for how
to evaluate them. It also has counterexamples: computer programs have generated conjectures for
decades~\cite{graffiti}, and machine learning has suggested new conjectures in knot theory and
representation theory~\cite{davies}. What remains open is whether machines will reframe
foundations in the way that the question ``what is geometry?'' once did.

\section{Conclusion}

Under a formalist definition, mathematics does not need a human audience, and a Thurston-style
definition becomes formalist as the class of understanders widens. Mathematics is the domain
where trust can most cleanly shift from producers to small, auditable checkers. Elsewhere,
trust in a superhuman agent can rationally exceed trust in human experts after finitely many
verified successes, for any standard short of certainty. Where the threshold actually sits is a
social fact, not a theorem. The decisive open assumption is whether reliability observed on
checkable tasks transfers to uncheckable ones.

\section{Lean index}\label{sec:lean}

\begin{center}
\small
\begin{tabular}{@{}rl@{}}
\toprule
Thm & Lean name (namespace \lean{TrustThesis}) \\
\midrule
1 & \lean{formalMath_producer_irrelevant} \\
2 & \lean{formalMath_of_humanistMath} \\
3 & \lean{humanistMath_mono} \\
4 & \lean{humanistMath_univ_iff} \\
5 & \lean{humanistMath_depends_on_boundary} \\
6 & \lean{trusted_base} \\
7, 8 & \lean{audit_miss_le}, \lean{audit_risk_eventually_lt} \\
9, 10 & \lean{posterior_tendsto_one}, \lean{trust_threshold_crossed} \\
11 & \lean{posterior_lt_one} \\
12 & \lean{outweighs_community} \\
13, 14 & \lean{evidence_blind}, \lean{exists_indistinguishable_failure} \\
15 & \lean{replacement_dilemma} (via \lean{sorites}) \\
16 & \lean{duplicate_understands} \\
\bottomrule
\end{tabular}
\end{center}

\begin{thebibliography}{99}
\bibitem{appelhaken} K.~Appel and W.~Haken, Every planar map is four colorable,
\emph{Illinois J.\ Math.} 21 (1977).
\bibitem{bpt} M.~J.~H.~Heule, O.~Kullmann and V.~W.~Marek, Solving and verifying the Boolean
Pythagorean triples problem via cube-and-conquer, \emph{SAT 2016}, LNCS 9710.
\bibitem{chalmers} D.~J.~Chalmers, Absent qualia, fading qualia, dancing qualia, in
T.~Metzinger (ed.), \emph{Conscious Experience}, 1995.
\bibitem{condorcet} Marquis de Condorcet, \emph{Essai sur l'application de l'analyse \`a la
probabilit\'e des d\'ecisions rendues \`a la pluralit\'e des voix}, 1785.
\bibitem{davies} A.~Davies et al., Advancing mathematics by guiding human intuition with AI,
\emph{Nature} 600 (2021), 70--74.
\bibitem{dietvorst} B.~J.~Dietvorst, J.~P.~Simmons and C.~Massey, Algorithm aversion: people
erroneously avoid algorithms after seeing them err, \emph{J.\ Exp.\ Psychol.\ Gen.} 144 (2015),
114--126.
\bibitem{graffiti} S.~Fajtlowicz, On conjectures of Graffiti, \emph{Discrete Math.} 72 (1988),
113--118.
\bibitem{flyspeck} T.~Hales et al., A formal proof of the Kepler conjecture, \emph{Forum Math.\
Pi} 5 (2017), e2.
\bibitem{debate} G.~Irving, P.~Christiano and D.~Amodei, AI safety via debate,
arXiv:1805.00899 (2018).
\bibitem{keller} J.~Brakensiek, M.~Heule, J.~Mackey and D.~Narv\'aez, The resolution of
Keller's conjecture, \emph{IJCAR 2020}.
\bibitem{lean4} L.~de Moura and S.~Ullrich, The Lean 4 theorem prover and programming
language, \emph{CADE-28} (2021).
\bibitem{paulson} L.~C.~Paulson, A machine-assisted proof of G\"odel's incompleteness theorems
for the theory of hereditarily finite sets, \emph{Rev.\ Symb.\ Logic} 7 (2014).
\bibitem{thurston} W.~P.~Thurston, On proof and progress in mathematics, \emph{Bull.\ Amer.\
Math.\ Soc.} 30 (1994), 161--177.
\bibitem{zach} R.~Zach, Hilbert's program, \emph{Stanford Encyclopedia of Philosophy}.
\end{thebibliography}

\end{document}
